% !PRINTED 444
\weberpage{444}
%<<CONT
\gesperrt{1. There are 24 sevenfold poles, which, corresponding to the eight seven-term cycles, divide into eight triples. Each of these 24 points can be transformed into every other by substitutions of the group $G_{168}$.}

\begin{center}
\S.\,118.

\gesperrt{The Principal Axes.}
\end{center}

We pass on to the consideration of the substitutions of degrees $4$ and $2$, as whose representatives we choose $\Theta$ and $\Theta^2$. According to \S.\,115, (10) and \S.\,116 we have to solve for $\Theta$ the cubic equation:
\weq{1}{
\begin{vmatrix}
\gamma\varepsilon^2-\lambda, & \alpha\varepsilon^6, & \beta\\
\alpha, & \beta\varepsilon^4-\lambda, & \gamma\varepsilon^5\\
\beta\varepsilon^3, & \gamma, & \alpha\varepsilon-\lambda
\end{vmatrix} = 0,
}
which, developed, reads:
\weq{2}{
\begin{aligned}
&\lambda^3-\lambda^2(\alpha\varepsilon+\beta\varepsilon^4+\gamma\varepsilon^2)\\
&\quad -\lambda\left[\varepsilon^6(\alpha^2-\beta\gamma)+\varepsilon^3(\beta^2-\gamma\alpha)+\varepsilon^5(\gamma^2-\alpha\beta)\right]\\
&\quad +\alpha^3+\beta^3+\gamma^3-3\alpha\beta\gamma=0.
\end{aligned}
}
But this equation reduces, with the use of the relations (8), (9), (14) of \S.\,115, to
\weq{3}{\lambda^3-\lambda^2+\lambda-1=0,}
and has the roots
\[
\lambda=1,\quad \lambda=\pm i,
\]
and from these the coordinates of the three poles can be computed, which result as rational functions of $\varepsilon$ and $i$.

For the poles of $\Theta^2$ we obtain, according to \S.\,115, (16), first the cubic equation
\[
\begin{vmatrix}
\beta-\lambda, & \gamma\varepsilon^3, & \alpha\varepsilon^2\\
\gamma\varepsilon^{-3}, & \alpha-\lambda, & \beta\varepsilon^{-1}\\
\alpha\varepsilon^{-2}, & \beta\varepsilon, & \gamma-\lambda
\end{vmatrix} = 0,
\]
which, by the relations between $\alpha,\beta,\gamma$, is easily brought into the form
\weq{4}{(\lambda+1)^2(\lambda-1)=0}
and therefore has a double root $\lambda=-1$ and a simple root $\lambda=1$. If we set the value $\lambda=-1$
%CONT>>
